Li et al.~\cite{li2026diverseevidence} study binary-event forecasting with a fixed evidence pool, partitioned into shared public and disjoint private subsets. Their InfoDelphi framework exchanges rationales during deliberation and combines final forecasts by confidence-weighted logit pooling; its outcomes are aggregate forecast accuracy and Brier score. They also evaluate a no-deliberation variant, so fixed evidence and independent outputs are not novel by themselves. Yang et al.~\cite{yang2026agentscaling} explain diminishing returns from correlated agent outputs through effective information channels and evaluate voting and debate systems. Their label-free proxy $K^*$ is the entropy effective rank of a Gram matrix of output embeddings, measuring semantic diversity rather than portfolio positions.

We share the principle that nominal agent count need not equal effective diversity, but target a different estimand. With the requested model roster, mandate, output rule, and team-level point-in-time evidence union fixed, we vary field access and retain each continuous signed book before communication or aggregation. We measure dependence through the cross-sectional book-correlation spectrum, using participation-ratio effective dimension, and complement it with same-side overlap and the proposed order footprint after internal netting. The contrast is therefore access-induced breadth of unaggregated positions, not deliberation gains, aggregate forecast quality, or semantic diversity: different reasoning can induce the same trade, while opposite books share a risk direction yet offset proposed demand.
